% Generated draft: requires mathematical and source review.
\documentclass[11pt]{article}
\usepackage[margin=25mm]{geometry}
\usepackage{fontspec}
\setmainfont{Latin Modern Roman}
\usepackage{amsmath,amssymb,mathtools,unicode-math}
\setmathfont{Latin Modern Math}
\usepackage{enumitem,needspace,xurl,hyperref}
\setlength{\parindent}{0pt}
\setlength{\parskip}{5pt}
\setlength{\emergencystretch}{3em}
\newcommand{\N}{\mathbb{N}}
\newcommand{\Z}{\mathbb{Z}}
\newcommand{\Q}{\mathbb{Q}}
\newcommand{\R}{\mathbb{R}}
\newcommand{\C}{\mathbb{C}}
\newcommand{\F}{\mathbb{F}}
\newcommand{\A}{\mathbb{A}}
\newcommand{\PP}{\mathbb{P}}
\newcommand{\E}{\mathbb{E}}
\newcommand{\eps}{\varepsilon}
\newcommand{\dd}{\,\mathrm{d}}
\begin{document}
\section*{UnsolvedMath ID: 30002345}
\textbf{Source problem code:} catalogue entry
\subsection{Definitions and mathematical statement}
The supplied question asks for which subfields \(K\subseteq\C\) the graded algebras of \(K\)-rational Hodge classes are isomorphic for every pair of conjugate varieties. The record does not specify the class of varieties or the precise meaning of algebra isomorphism. The following formulation is conditional on interpreting the varieties as smooth projective complex varieties and the isomorphisms as degree-preserving \(K\)-algebra isomorphisms; these conventions require source confirmation.

Under these conventions, let \(X\) be a smooth projective complex variety. Using extension of coefficients to regard singular cohomology \(H^{2p}(X(\C),K)\) as a subspace of \(H^{2p}(X(\C),\C)\), define
\[
H^{p,p}(X,K):=H^{2p}(X(\C),K)\cap H^{p,p}(X),\qquad
H^{*,*}(X,K):=\bigoplus_{p=0}^{\dim X}H^{p,p}(X,K).
\]
Here \(H^{p,p}(X)\) is the \((p,p)\)-summand of the complex Hodge decomposition. Equip this algebra with the cup product and give its \(p\)-th summand cohomological degree \(2p\). For \(\sigma\in\operatorname{Aut}(\C)\), let \(X^\sigma\) denote the complex variety obtained by applying \(\sigma\) to the coefficients of defining equations for \(X\).

With the stated conventions, determine the set
\[
\left\{K\subseteq\C\ \middle|\ 
\begin{gathered}
K\text{ is a subfield, and for every smooth projective complex variety }X\\
\text{and every }\sigma\in\operatorname{Aut}(\C),\text{ there exists a degree-preserving}\\
K\text{-algebra isomorphism }H^{*,*}(X,K)\cong H^{*,*}(X^\sigma,K)
\end{gathered}
\right\}.
\]
No canonical isomorphism or compatibility among these isomorphisms is requested. This conditional formulation places no restriction \(\sigma(K)=K\); whether the source intends such a restriction remains unconfirmed. The supplied record labels the problem partially solved, but that label and its associated literature assessment are supplied curation rather than independently verified mathematical conclusions.
\subsection{Short English statement}
Determine which subfields \(K\subseteq\C\) make the graded algebra of \(K\)-valued Hodge classes invariant, up to isomorphism, under conjugation of every variety in the intended class. The displayed formulation uses smooth projective complex varieties and degree-preserving \(K\)-algebra isomorphisms as explicit provisional conventions; the supplied record leaves these details unspecified.
\subsection{Sources}
\begin{itemize}
\item Local UnsolvedMath registry, supplied complete record: permanent ID \(30002345\), title \emph{Hodge Algebras of Conjugate Varieties}, exact source problem code \texttt{catalogue entry}, Oberwolfach Reports Open Problems collection. The record gives proposed year 2013 and names no proposer.
\item Supplied report citation: \emph{Complex Algebraic Geometry (2014)}, Oberwolfach Reports, pp.~1563--1627, DOI \texttt{10.4171/owr/2013/27}. \url{https://doi.org/10.4171/owr/2013/27}. The supplied title year is 2014, whereas the proposed year and DOI contain 2013; this discrepancy is unresolved.
\item Schreieder, \emph{Multiplicative sub-Hodge structures of conjugate varieties}. Supplied literature reference: \url{https://arxiv.org/abs/1304.5146}.
\item Supplied \emph{Forum of Mathematics Sigma} publication reference: \url{https://doi.org/10.1017/fms.2013.7}.
\item Literature assessment dated exactly \texttt{2026-08-21}: the supplied curation labels the problem partially solved and attributes to Schreieder the assertion that fields not contained in an imaginary quadratic field admit conjugate varieties with different rational \((p,p)\)-class algebras; it further states that the remaining imaginary-quadratic cases reduce to the Hodge conjecture and are unresolved. This is supplied dated curation, not an independently verified assessment. The background expressly requests independent verification of time-sensitive status and cited claims.
\item The supplied registry marks the statement \texttt{corrected\_verified} and reports a manual semantic, typographic, and LaTeX review of its cleanup. This records the supplied editorial assessment and does not establish the missing mathematical conventions or constitute independent source verification.
\end{itemize}
\end{document}
